% Transcription — fonds Alexandre Grothendieck, shelfmark 161-2, batch 1 (pages 1-20).
% Established from the facsimile archives/batches/161-2/batch-01.pdf (a mirror of
% https://grothendieck.umontpellier.fr/161-2.pdf), per the skill
% transcribe-grothendieck. Page 1 is the folder's cover, whose tab supplies the
% opening heading; pages 2 and 16 are blank — all three absent.
% Pass: Opus 5 (claude-opus-5), 2026-09-02 — first pass, unchecked against the pages by a human.
\documentclass[11pt,a4paper]{article}
% Le préambule commun à toutes les transcriptions du fonds Grothendieck.
%
% Il tient en une idée : **l'incertitude se compose, elle ne se résout pas.**
% Une transcription qui lisse un mot douteux a détruit la seule chose pour
% laquelle on la faisait — un lecteur doit pouvoir savoir, à chaque endroit,
% ce qui était sur le feuillet et ce qui a été ajouté. Les six macros qui
% suivent sont donc l'appareil critique, et non de la décoration.
%
% Les mêmes commandes sont reconnues par `scripts/render.mjs`, qui produit la
% vue de lecture du site. Une macro ajoutée ici doit l'être aussi là-bas,
% faute de quoi le rendu échouera bruyamment — ce qui est voulu : un rendu
% silencieusement faux est pire qu'un rendu absent.

\ProvidesPackage{grothendieck}[2026/08/08 Transcriptions du fonds Grothendieck]

\usepackage{amsmath,amssymb,amsthm}
\usepackage{mathrsfs}
\usepackage{stmaryrd}
% Les diagrammes commutatifs sont partout dans ce fonds : c'est la langue
% même de la géométrie algébrique de Grothendieck, et une transcription qui
% les rendrait en prose ne transcrirait rien.
\usepackage{tikz-cd}
% Français par défaut — c'est la langue des feuillets — mais l'anglais est
% chargé aussi : la traduction et le résumé sont des documents anglais, et la
% typographie française (espaces avant les deux-points) y serait une faute.
% Ces documents basculent par \selectlanguage{english} en tête de corps.
\usepackage[english,french]{babel}
\usepackage{fontspec}
\usepackage{geometry}
\geometry{a4paper,margin=25mm,marginparwidth=28mm}
\usepackage{marginnote}
\usepackage{xcolor}
% ulem, not soul. `soul`'s \st and \ul are built on a character-by-character
% scan that XeLaTeX's font machinery breaks: the document fails with an
% undefined control sequence rather than degrading. ulem does the same two jobs
% and is Unicode-engine safe. `normalem` stops it hijacking \emph, which the
% transcriptions use for Grothendieck's own emphasis.
\usepackage[normalem]{ulem}
\usepackage{hyperref}
\hypersetup{colorlinks=true,linkcolor=black,urlcolor=black,pdfborder={0 0 0}}

% --- Métadonnées du lot -----------------------------------------------------
% Reprises telles quelles par le script de rendu, qui les affiche en tête de
% la vue de lecture. Elles fixent ce que le fichier prétend couvrir : sans
% elles, une transcription flotte sans point d'attache dans un fonds de seize
% mille feuillets.

\newcommand{\folder}[1]{\gdef\grothFolder{#1}}
\newcommand{\batch}[1]{\gdef\grothBatch{#1}}
\newcommand{\leaves}[2]{\gdef\grothFirst{#1}\gdef\grothLast{#2}}
\newcommand{\foldertitle}[1]{\gdef\grothFoldertitle{#1}}
\gdef\grothFolder{?}\gdef\grothBatch{?}\gdef\grothFirst{?}\gdef\grothLast{?}\gdef\grothFoldertitle{}

% --- Le résumé --------------------------------------------------------------
%
% Il ouvre la lecture modernisée, sous le titre « Résumé », et il est composé
% en petit. Ce n'est pas un ornement : c'est le seul passage du document qui
% ne soit pas des mathématiques, et un lecteur doit pouvoir le reconnaître
% avant de l'avoir lu — puis le sauter s'il connaît déjà le sujet, ou s'y
% arrêter s'il ne le connaît pas. Le filet qui le suit marque la reprise du
% corps.

\newenvironment{resume}
  {\small\setlength{\parskip}{0.45em}}
  {\par\medskip\hrule\medskip}

% --- Mots-clés ---------------------------------------------------------------
%
% En anglais, en clôture du résumé de la lecture modernisée : le vocabulaire
% moderne sous lequel chercher ce que les feuillets construisent. Le manifeste
% du site (`npm run manifest`) les extrait pour étiqueter la cote entière —
% cette ligne est la source unique des tags, il n'y a pas de fichier à côté.
\newcommand{\keywords}[1]{\par\smallskip\noindent
  {\footnotesize\textsc{Keywords} --- \textit{#1}}\par}

% --- L'appareil critique ----------------------------------------------------

% Le numéro de feuillet, en marge. C'est l'ancre qui permet de régler un
% désaccord sur une formule en regardant *un* feuillet plutôt qu'une chemise
% de six cents. Le site s'en sert aussi pour tourner les pages du fac-similé
% quand on fait défiler la transcription.
\newcommand{\leaf}[1]{%
  \marginnote{\footnotesize\color{gray!70}\textsf{#1}}%
  \label{leaf:#1}\ignorespaces}

% Illisible : on ne devine pas. Un mot inventé qui se lit comme les autres est
% le pire résultat possible.
\newcommand{\ill}{\textcolor{red!60!black}{[\,\ldots\,]}}

% Lecture douteuse : le mot est donné, et signalé comme incertain.
\newcommand{\uncertain}[1]{\uline{#1}}

% Ajout de l'éditeur — une lettre manquante, un mot sous-entendu.
\newcommand{\add}[1]{\textcolor{blue!50!black}{[#1]}}

% Biffé par Grothendieck lui-même. Ce qu'il a rayé fait partie du document :
% une rature dit souvent où la pensée a changé de direction.
\newcommand{\struck}[1]{\sout{#1}}

% Note du transcripteur, sur l'état matériel du feuillet.
\newcommand{\note}[1]{\par\smallskip{\small\color{gray!60!black}\textit{[#1]}}\par\smallskip}

% Note marginale de Grothendieck, distincte du corps du texte.
\newcommand{\marginal}[1]{\par\smallskip\noindent\hspace{1em}%
  {\small\color{gray!60!black}#1}\par\smallskip}

% --- Titre ------------------------------------------------------------------

\AtBeginDocument{%
  \begin{center}
    {\large\bfseries Fonds Alexandre Grothendieck --- cote n\textsuperscript{o}~\grothFolder}\\[2pt]
    {\itshape\grothFoldertitle}\\[4pt]
    {\small feuillets \grothFirst--\grothLast\ (lot \grothBatch)}\\[2pt]
    {\footnotesize\color{gray!60!black}Transcription établie d'après le fac-similé numérisé
     par l'Université de Montpellier.\\
     Les lectures incertaines sont soulignées, les illisibles notées [\,\ldots\,],
     les ajouts entre crochets.}
  \end{center}
  \vspace{1em}}

\endinput


\folder{161-2}
\batch{1}
\pages{1}{20}
\dating{[vers 1963-1973]}
\watermark{Édition de démonstration}
\foldertitle{Algèbre universelle [ou catégories] : notes manuscrites (s.d.)}

\begin{document}

\section*{Algèbre universelle}

\note{titre porté de sa main sur l'onglet de la chemise (feuillet 1, non transcrit)}

\page{3}
\textbf{Théorème}\quad Soient $d = (d', d'')$ un \add{couple de} petits ens. de
diagrammes, $C$ une catégorie ess. petite, $\delta = (\delta', \delta'')$ un
\add{petit} ens. de cônes projectifs et inductifs dans $C$,
$\Delta = \Delta_{d} = \Delta_{d',d''}$ : la catégorie des catégories ayant des
$\lim$ de type $d$, avec foncteurs y commutant.

\struck{Alors il existe}

Pour tout $E \in \mathrm{Ob}\,\Delta$, soit
$\underline{\mathrm{Hom}}_{\Delta,\delta}(C,E)$ la sous-cat. pleine de
$\underline{\mathrm{Hom}}(C,E)$ formée des foncteurs qui sont $\delta$-exacts.
Si \ill{} est de même de $(\delta'_1, \delta''_1)$ et est $\in \Delta$, on
définit de même $\underline{\mathrm{Hom}}_{\Delta,\delta_1}(C_1,E)$. Ceci posé,
il existe une catégorie $\bar{C}_{\delta} \in \mathrm{Ob}\,\Delta$, et un
foncteur $i : C \to \bar{C}_{\delta}$, tel que

\begin{enumerate}
  \item[a)] $i$ soit $\delta$-exact, i.e.
    $i \in \underline{\mathrm{Hom}}_{\delta}(C, \bar{C}_{\delta})$ ;
  \item[b)] il soit universel en le sens suivant : pour tout
    $E \in \mathrm{Ob}\,\Delta$, le foncteur $\varphi \mapsto \varphi i$ induit
    une équivalence
    \[
      \underline{\mathrm{Hom}}_{\Delta,\, i(\delta)}(\bar{C}_{\delta}, E)
      \;\xrightarrow{\ \approx\ }\;
      \underline{\mathrm{Hom}}_{\delta}(C, E)
    \]
\end{enumerate}

[i.e. $(\bar{C}_{\delta}, i)$ est objet $2$-initial de la $2$-catégorie des
\struck{topos} $\in \mathrm{Ob}\,\Delta$, \uncertain{$E$} tels que $C \to E$ soit
$\delta$-exact]. De plus, $\bar{C}_{\delta}$ est ess. petite.
\note{le glyphe lu ici $E$ a la forme du $C$ qu'il trace ailleurs ; le sens exige un objet de $\Delta$}

\textbf{N.B.}\quad Ce Pb redonne à la fois le Pb 1 ($\delta' = \delta'' =
\emptyset$) et le Pb 2 ($C \in \mathrm{Ob}\,\Delta$, \struck{\ill{}}
$\delta' = $ \struck{diagrammes} \add{cônes pr. de type $d'$} exacts
\struck{égaux} $\cup$ les diagrammes $(x \to y) \in \Sigma$,
$\delta'' = $ \struck{diagrammes} cônes inductifs de type $d''$ exacts). On verra
que le Th. permet aussi de traiter les cas de conditions d'exactitude.

\textbf{Dém.}\quad On construit \add{par} \struck{\ill{}} récurrence transfinie
\add{une} suite transfinie de \add{petites} catégories $C_i$ sous $C$, munies de
$\delta_i = (\delta'_i, \delta''_i)$ (petit) :
\[
  C \;\xrightarrow{\ \alpha_i\ }\; C_i \;\xrightarrow{\ \tau_i\ }\; C_{i+1}
\]
\[
\begin{cases}
  C_0 = C, \quad \delta_0 = \delta = (\delta', \delta''), \\
  (C_{i+1}, \delta_{i+1}) = K(C_i, \delta_i), \\
  C_i = \varinjlim_{j<i,\ (\mathrm{Cat})} C_j \quad \text{si $i$ ordinal limite.}
\end{cases}
\]
\struck{$(C_{i+1}, \delta'_{i+1})$ déduit de $(C_i, \delta_i)$ ainsi :} — $K$
étant explicité ci-dessous ; et, pour $i$ ordinal limite,
$\delta'_i = \bigcup \mathrm{Im.}\ \delta'_j$ dans $C_i$,
$\delta''_i = \bigcup \mathrm{Im.}\ \delta''_j$ dans $C_i$.

\page{4}
Les propriétés essentielles de $K(C,\delta)$ sont résumées dans :

\begin{enumerate}
  \item[a)] $\tau_i : C_i \to C_{i+1}$ est \struck{préexact} \add{$\delta_i$-}préexact$_{\delta}$, i.e.
  \item[$a'_1$)] \struck{Une double flèche $z \rightrightarrows x$} Pour
    $(x \xrightarrow{\,p_\lambda\,} x_\lambda) \in d'_i$, et tout \struck{dans le}
    couple de flèches $z \overset{u}{\underset{v}{\rightrightarrows}} x$ telles
    que $p_\lambda u = p_\lambda v$ $\forall \lambda$, on a
    $\tau_i(u) = \tau_i(v)$.
    \note{il écrit ici $d'_i$ ; plus bas, feuillet 5, la construction $\rho$ porte sur $\delta''$}
  \item[$a'_2$)] Pour toute \struck{\ill{}} augmentation
    $(z \xrightarrow{\,q_\lambda\,} x_\lambda)$ du syst. proj., il existe
    $\tau_i(z) \dashrightarrow \tau_i(x)$ telle que
    $\tau_i(p_\lambda) \circ q = \tau_i(q_\lambda)$ $\forall \lambda$.
  \item[$a''$)] \struck{\ill{}} les mêmes, avec $d''_i$ au lieu de $d'_i$.
\end{enumerate}

\struck{b) Tout diagramme $d'_{i+1} \supset d' \to \delta_{i+1} \to \tau_i(\delta_i)$}

\begin{enumerate}
  \item[$b'$)] $\forall$ diagramme \add{$\Theta'$} de type $d'$ dans $C_i$,
    $\exists$ une \struck{\ill{}} augmentation de $\tau_i(\Theta')$ qui soit
    $\in \delta'_{i+1}$ ;
  \item[$b''$)] $\forall$ diagramme \add{$\Theta''$} de type $d''$ dans $C_i$,
    $\exists$ une augmentation de $\tau_i(\Theta'')$ \add{dans $\delta''_{i+1}$}.
  \item[c)] $\delta'_{i+1} = $ \struck{\ill{}} $\tau_i(\delta'_i) \cup$
    diagrammes de type $d'$ \struck{\ill{}} ;
    $\delta''_{i+1} = \tau_i(\delta''_i) \cup$ diagrammes de type $d''$.
    \struck{\ill{} de $\delta'_{i+1}$ de type $\Theta'$ \ill{} et $d'$
    ($\Theta'' \in d''$) est dans $\tau_i(\delta'_i)$, $(\tau_i(\delta''_i))$}
  \item[d)] $\forall$ catégorie $E \in \mathrm{Ob}\,\Delta$, le foncteur
    $\varphi \mapsto \varphi \circ \tau_i$ induit une équivalence
    \[
      \underline{\mathrm{Hom}}_{\delta_{i+1}}(C_{i+1}, E)
      \;\longrightarrow\;
      \underline{\mathrm{Hom}}_{\delta_i}(C_i, E).
    \]
\end{enumerate}

\struck{La} La construction de $K : (C,\delta) \mapsto K(C,\delta)$ sera donnée
plus bas. On \struck{peut} \uncertain{voit} maintenant comment elle permet de
terminer.

\marginal{de c) on conclut par récurrence transfinie : $(c^{*})$
$\delta'_i = \alpha_i(\delta') \cup$ \struck{\ill{}} \add{images} des diagrammes
de type $d'$ \struck{\ill{}} ; $\delta''_i = \alpha_i(\delta'') \cup$
\struck{\ill{}} \ill{} de type $d''$.}

\struck{De c) on conclut par récurrence transfinie}

De d), on déduit par récurrence transfinie
\[
  (d^{*}) \qquad
  \boxed{\ \underline{\mathrm{Hom}}_{\delta_i}(C_i, E)
  \;\xrightarrow{\ \approx\ }\;
  \underline{\mathrm{Hom}}_{\delta}(C, E)\ }
  \qquad \varphi \mapsto \varphi \circ \alpha_i .
\]
\marginal{ordinal ?}

Donc on a terminé si on \uncertain{montre} qu'il existe $\omega$ tel que

\begin{enumerate}
  \item[$1^{\circ}$)] $C_\omega \in \mathrm{Ob}\,\Delta$ ;
  \item[$2^{\circ}$)] $\alpha_\omega : C \to C_\omega$ est $\delta$-exact ;
  \item[$3^{\circ}$)] $\underline{\mathrm{Hom}}_{\delta_\omega}(C_\omega, E)
    = \underline{\mathrm{Hom}}_{\Delta,\, \uncertain{\alpha_\omega}(\delta)}(C_0, E)$.
\end{enumerate}

\struck{\ill{}} Soit pour ceci $\pi$ un cardinal tel que
$\pi > \mathrm{card}\,\mathrm{Fl}\,\Theta$, \struck{\ill{}} si $\Theta \in d'$ ou
$\Theta \in d''$, ou $\Theta$ \struck{\ill{}} indexe un $\in \delta'$ ou un
$\in \delta''$ ; et $\omega$ le plus petit ordinal de cardinal $\geqslant \pi$.
Donc \ill{}

\page{5}
$j < \omega$ est grand devant $\mathrm{card}\,\mathrm{Fl}\,\Theta$, pour tout
$\Theta$ comme dessus.

On va établir d'abord le point suivant, plus fort que $2^{\circ}$) :
\[
  (*) \qquad \forall\, j < \omega, \quad C_j \to C_\omega \ \text{est}\ \delta_j\text{-exact}.
\]
Cela résulte aisément de a) et c) \add{et de la grandeur de $\omega$}.
\marginal{Détailler la démonstration}

On \struck{peut} \add{en conclut} $1^{\circ}$), utilisant $(*)$ et b).

Pour prouver $3^{\circ}$), on note que b) \add{et c)} et la grandeur de $\omega$
indiquent
\[
  (**) \qquad \delta_\omega = \text{(cônes exacts de types $d'$ et $d''$)}
  \ \cup\ \alpha_\omega(\delta).
\]

Tout revient donc à construire $K(C,\delta)$, de façon à satisfaire a) ; d).

On construit \add{$(C,\delta) \mapsto$} $K(C,\delta)$ comme composé
$(C,\delta) \mapsto \sigma\rho(C,\delta)$ de \uncertain{deux} constructions,
dont chacune fait environ la moitié du boulot :

$(C,\delta) \mapsto \rho(C,\delta)$ se fait en plusieurs étapes. On considère
$\Delta_{(\emptyset,\, d'')} = \Delta''$, $\Delta_{(d',\, \emptyset)} = \Delta'$.
\marginal{Dans chacune, on doit vérifier l'analogue des c) et d) (ce qui sera vrai pour la composée)}

\struck{plus précis :}

$\rho_1(C,\delta) = (C_1, \delta_1)$, où $C_1$ est la $\Delta''$-enveloppe de
$C$, $\tau_0 : C_0 \to C_1$ ;
$\delta'_1 = \tau_0(\delta')$,
$\delta''_1 = \tau_0(\delta'') \cup ($diagrammes de type $d''$ de $C_1)$.
\marginal{Cela assure \uncertain{b')}\quad c) et d) vérifiés}

$\rho_2(C_1, \delta_1) = (C_2, \delta_2)$, \add{$\tau_1 : C_1 \to C_2$}, avec
$C_2$ déduit de $C_1$ en égalisant toutes les couples de flèches
$z \overset{u}{\underset{v}{\rightrightarrows}} x$ dans $C_1$, telles que
\ill{} $x$ d'une $(x \xrightarrow{\,p_\lambda\,} x_\lambda) \in \delta''_1$,
satisfaisant $p_\lambda u = p_\lambda v$ pour tout $\lambda$. On a
$\delta_2 = \tau_1(\delta_1)$.
\marginal{Cela assure $a'_1$.\quad c) et d) sont vérifiés}

$\rho_3(C_2, \delta_2) = (C_3, \delta_3)$, \add{$\tau_2 : C_2 \to C_3$}, avec
$C_3$ déduit de $C_2$ en rajoutant, pour tout
$\Theta'' = (x \xrightarrow{\,p_\lambda\,} x_\lambda) \in \delta''_2$ et toute
augmentation $q = (z \xrightarrow{\,q_\lambda\,} x_\lambda)_\lambda$ du système
projectif $(x_\lambda)$, une flèche
$u_{\Theta'',\, q} : z \to x$ dans $C_3$, satisfaisant
$p_\lambda u_{\Theta'',\, q} = q_\lambda$ $\forall \lambda$. On a
$\delta_3 = \tau_2(\delta_2)$.
\marginal{Cela assure $a'_2$.\quad c) et d) vérifiés}

\page{6}
\[
  \rho(C,\delta) = \rho_3 \rho_2 \rho_1 (C,\delta).
\]
On définit de façon analogue une \struck{opération} \add{construction}
$\sigma = \sigma_3 \sigma_2 \sigma_1$, assurant successivement $b''$),
\struck{$b''_1$} $a''_1$, $a''_2$. On pose
\[
  K(C,\delta) = \sigma\rho(C,\delta)
  = \sigma_3 \sigma_2 \sigma_1 \rho_3 \rho_2 \rho_1 (C,\delta).
\]
Il a les propriétés a) à d) voulues.

\textbf{N.B.}\quad L'ordre dans lequel on a \ill{} n'est pas d'importance. On
pourrait \ill{} prendre par ex.
\[
  K'(C,\delta) = \sigma_1 \rho_3 \sigma_2 \rho_1 \rho_2 \sigma_3 (C,\delta)\ !
\]
\ill{} composé les \add{$6$ opérations} $\sigma_i$, \struck{et les} $\rho_i$
\ill{}. La \struck{chose} \ill{} ayant une signification intrinsèque est la
limite des itérés transfinis de $K$ \ill{}

\section*{Conditions d'exactitude}

On se donne \add{des triplets} \struck{des petites catégories}
\add{des $\Delta$-foncteurs entre catégories} dans
$\Delta = \Delta_d = \Delta_{(d',d'')}$,
\[
  f^{(i)}_\lambda : R^{(i)}_\lambda \longrightarrow R'^{(i)}_\lambda
  \qquad (i = 0,1,2),
\]
\ill{} $f^{i} = (f_\lambda)$, $f = (f^{(0)}, f^{(1)}, f^{(2)})$.
\marginal{(petits)}

$\Delta_f$ = la sous-$2$-catégorie pleine de $\Delta$ formée des
$E \in \mathrm{Ob}\,\Delta$ tels que, \struck{le foncteur $\varphi \mapsto \varphi f$}
pour tout $\varphi : R \to R'$ dans $\Delta$, les foncteurs
$u \mapsto u \circ \varphi$, de
$\underline{\mathrm{Hom}}_{\Delta}(R', E) \to \underline{\mathrm{Hom}}_{\Delta}(R, E)$,
soient

\begin{enumerate}
  \item[$a_0$)] surjectifs si $\varphi$ est un $f^0_\lambda$ ;
  \item[$a_1$)] \struck{sont} bijectifs \struck{\ill{}} \add{injectifs} si $\varphi$ est un $f^1_\lambda$ ;
  \item[$a_2$)] \ill{} si $\varphi$ est un $f^2_\lambda$.
\end{enumerate}
\marginal{(Les foncteurs \add{sur $E$} \ill{} les foncteurs \ill{} exacts)}

Avec les notations précédentes :

\textbf{Théorème}\quad [Soit $(C,\delta)$ comme dans le Th. Alors on peut trouver
un foncteur \add{$\delta$-exact} $C \xrightarrow{\ i\ } \bar{C}$, avec
$\bar{C} \in \mathrm{Ob}\,\Delta_f$, et qui soit \add{$2$-}universel pour ces
propriétés. De plus, $\bar{C}$ est ess. petite. Cela revient en particulier
\ill{} Pb 1 et Pb 2 \add{et Pb $2'$} (pour $\Delta_f$).]

\page{7}
\struck{On peut faire la construction de $K(\bar{C}, \bar{\delta})$ en trois
étapes, par construction transfinie comme
$K(C,\delta) = \rho_2 \rho_1 \rho_0 (C,\delta)$, avec $\rho_0$) précédemment,
\ill{} sont des conditions, $\rho_2$) aux \ill{} a) b) c). $\rho_0$) Si
$R \xrightarrow{\varphi} R'$ est \ill{} $f^i_\lambda$, et
$u, v : R' \rightrightarrows C_i$.}
\note{tout ce début de feuillet est biffé de deux longs traits diagonaux}

Utilisant le Th. précédent, on peut supposer que $C \in \mathrm{Ob}\,\Delta$, et
$\delta'$ $(\delta'')$ \struck{contient les} \struck{diagrammes} \ill{} (exacts
de type $d'$ $(d'')$), de sorte que
\[
  \underline{\mathrm{Hom}}_{\delta}(C, E) = \underline{\mathrm{Hom}}_{\Delta,\delta}(C, E)
  \qquad \forall\, E \in \mathrm{Ob}\,\Delta .
\]
\marginal{(petite)}

On va faire une construction transfinie $(C_i, \delta_i)$,
$C_i \in \mathrm{Ob}\,\Delta$, foncteurs de transition dans $\Delta$ :
\[
\begin{cases}
  (C_0, \delta_0) = (C, \delta), \\
  (C_{i+1}, \delta_{i+1}) = K(C_i, \delta_i), \quad
    \tau_i : C_i \to C_{i+1}, \quad \delta_{i+1} = \tau_i(\delta_i), \\
  C_i = \varinjlim_{j < i,\ (\Delta)} C_j \ \text{ si $i$ ordinal limite,} \quad
    \delta_i = \textstyle\bigcup_{j<i} \mathrm{Im.\ de\ } \delta_j \ \text{dans } C_i .
\end{cases}
\]

On construit $(C,\delta) \mapsto K(C,\delta)$ de façon à satisfaire les
conditions suivantes :

\begin{enumerate}
  \item[$a_0$)] Pour tout foncteur $R \xrightarrow{\ \varphi\ } R'$ parmi les
    $f^{(i)}_\lambda$ $(i = 0,1,2)$, et pour tout \add{couple de} foncteurs
    $u', v' : R' \rightrightarrows C_i$, et une double flèche
    $\lambda, \mu : u' \to v'$, si $\lambda * \varphi = \mu * \varphi$, alors
    $\tau_i(\lambda) = \tau_i(\mu)$.
  \item[$a_1$)] Pour tout \struck{\ill{}} \add{$\varphi$} \struck{foncteur}
    $\varphi : R \to R'$ parmi les $f^1_\lambda$ et $f^2_\lambda$, et tout
    couple de \add{$2$-}foncteurs $u', v' : R' \to C_i$, et
    $\lambda' : u'\varphi \to v'\varphi$, $\exists$ \struck{\ill{}} considérant
    $\tau_i \circ u' \xrightarrow{\ \tau_i * \lambda'\ } \tau_i \circ v'$,
    \[
      \Bigl( R \ \overset{\tau_i \circ u' \varphi}
      {\underset{\tau_i \circ v' \varphi}{\rightrightarrows}}\ C_{i+1} \Bigr),
    \]
    $\exists\, \lambda : \tau_i \circ u' \varphi \to \tau_i \circ v' \varphi$ tel
    que $\tau_i * \lambda' = \lambda * \varphi$.
  \item[$a_2$)] Pour tout $\varphi : R \to R'$ parmi les $f^2_\lambda$, et tout
    $2$-foncteur $u : R \to C_i$, $\exists$ \ill{}
    $\tau_i \circ u : R \to C_{i+1}$, un $2$-foncteur $u' : R' \to C_{i+1}$ et un
    isomorphisme $\tau_i \circ u \simeq u' \circ \varphi$.
\end{enumerate}

\page{8}
\begin{enumerate}
  \item[b)] $\tau_i$ est $\delta_i$-exact.
  \item[c)] Pour tout $E \in \mathrm{Ob}\,\Delta_f$, le foncteur
    $\varphi \mapsto \varphi \circ \tau_i$, de
    $\underline{\mathrm{Hom}}_{\Delta}(C_{i+1}, E) \to \underline{\mathrm{Hom}}_{\Delta}(C_i, E)$,
    induit une équivalence
    \[
      \underline{\mathrm{Hom}}_{\Delta,\delta_{i+1}}(C_{i+1}, E)
      \;\xrightarrow{\ \approx\ }\;
      \underline{\mathrm{Hom}}_{\Delta,\delta_i}(C_i, E).
    \]
\end{enumerate}

\struck{On pose}

De c) on déduit par récurrence transfinie, pour $E \in \mathrm{Ob}\,\Delta_f$,
\[
  (c^{*}) \qquad
  \underline{\mathrm{Hom}}_{\Delta,\delta_i}(C_i, E)
  \;\xrightarrow{\ \approx\ }\;
  \underline{\mathrm{Hom}}_{\Delta,\delta}(C, E),
\]
et de la définition des $\delta_i$, on déduit par réc. transfinie
$\delta_i = \alpha_i(\delta)$ $(\alpha_i : C \to C_i)$.

Donc tout revient à prouver qu'\struck{on obtient} il existe un $i$ \add{limite}
tel que \add{(petit)} $C_\omega \in \mathrm{Ob}\,\Delta_f$, et $C \to C_\omega$
soit $\delta$-exact. Pour ceci, soit $\pi$ un cardinal majorant strictement
$\mathrm{Fl}\,R$ et $\mathrm{Fl}\,R'$, \struck{$\mathrm{Fl}\,\Theta$} pour les
$R$, $R'$ qui interviennent dans $f$, ainsi que les $\mathrm{Fl}\,\Theta$ où
$\Theta$ intervient dans $\delta'$ ou $\delta''$ ; soit $\omega$ le plus petit
ordinal de cardinal $\pi$. Donc \ill{} $\delta_\omega$ \ill{} $j < \omega$
\ill{} est grand devant $\mathrm{card}\,\mathrm{Fl}\,R$,
$\mathrm{card}\,\mathrm{Fl}\,R'$, $\mathrm{Fl}\,\Theta$.

\marginal{\ill{} déjà sont \ill{} indépendants de catégories \ill{}}

Mais on voit que $C \to C_\omega$ est $\delta$-exact en utilisant
\add{card} $\omega$ grand devant les $\mathrm{card}\,\mathrm{Fl}\,\Theta$. Et on
\ill{} b) et le fait que $\delta_\omega$ soit \ill{} est que pour
$\varphi : R \to R'$,
$\underline{\mathrm{Hom}}(R', C_\omega) \to \underline{\mathrm{Hom}}(R, C_\omega)$
est surj. \struck{\ill{}} (\ill{} $f^{(i)}_\lambda$ \ill{} $(i = 0,1,2)$)
\ill{} pleinement fidèle $(i = 1,2)$ \ill{}, en utilisant respectivement
$a_0$, $a_1$, $a_2$ plus le fait que $\omega$ est grand devant
\struck{\ill{}} $\mathrm{card}\,\mathrm{Fl}\,R$ et
$\mathrm{card}\,\mathrm{Fl}\,R'$.

\page{9}
Il reste donc à faire la construction de $K(C,\delta)$.

\page{10}
\section*{Conditions d'exactitude}
\note{il reprend le titre du feuillet 6, et recommence}

\begin{enumerate}
  \item[$1$)] Sur catégories $C$.
  \item[a)] Existence inconditionnelle \add{dans $C$} de certaines
    $\varprojlim$ ou $\varinjlim$, correspondant à certains types de diagrammes
    $d = (d', d'')$.
  \item[b)] Existence conditionnelle des types de $\varprojlim$ ou
    $\varinjlim$. On se donne \struck{\ill{}} \add{un} diagramme $\Theta$, et
    \struck{des sous-\ill{}} \ill{}
    $\Theta'_i \xrightarrow{\ u'_i\ } \Theta$,
    $\Theta''_j \xrightarrow{\ u''_j\ } \Theta$, et on demande que pour tout
    $\Theta \to C$ (satisfaisant éventuellement des conditions supplémentaires :
    commutativité de certains diagrammes, si on remplace les $\Theta \to$ \ill{},
    ce qui revient à demander que $\Theta$ soit une catégorie, et
    $\Theta \to C$ un foncteur) [\ill{} il existe des conditions
    « d'exactitude » plus subtiles, cf. plus bas] les $\varprojlim$
    ($\varinjlim$) des foncteurs correspondants aux types $\Theta'_i$
    ($\Theta''_j$) existent dans $C$.
  \item[c)] Exactitude de cônes projectifs \struck{diagrammes} \ill{} et
    inductifs. On se donne \ill{} un paquet
    $(\delta', \lambda', \delta'', \lambda'')$ de cônes projectifs et inductifs,
    et on demande que tout foncteur $\delta = (\delta',\delta'')$-exact soit
    $\lambda = (\lambda', \lambda'')$-exact. C'est un type d'exactitude \ill{}
  \item[d)] On se donne un \add{foncteur} \struck{catégorie}
    $\Theta \xrightarrow{\ \varphi\ } \bar\Theta$ entre catégories, et dans
    $\Theta$ un paquet $\delta = (\delta', \delta'')$, dans $\bar\Theta$ un
    paquet $\bar\delta = (\bar\delta', \bar\delta'')$. On considère
    \add{$u \mapsto u \circ \varphi$}
\end{enumerate}

\begin{tikzcd}[column sep=small, row sep=small, nodes={font=\scriptsize}]
\underline{\mathrm{Hom}}(\bar\Theta, C) \arrow[r, "\varphi^{*}"] & \underline{\mathrm{Hom}}(\Theta, C) \\
\underline{\mathrm{Hom}}_{\bar\delta}(\bar\Theta, C) \arrow[u, hook] & \underline{\mathrm{Hom}}_{\delta}(\Theta, C) \arrow[u, hook]
\end{tikzcd}

\marginal{Exiger \ill{} l'existence \ill{} ?}

\page{11}
et on peut demander des conditions du type suivant :

\begin{enumerate}
  \item[(i)] $\varphi^{*}_C$ envoie $\underline{\mathrm{Hom}}_{\bar\delta}$ dans
    $\underline{\mathrm{Hom}}_{\delta}$ (on particularise \ill{} $C$) ;
  \item[(i$'$)] \struck{\ill{}} la restriction de $\varphi^{*}_C$ :
    $\underline{\mathrm{Hom}}_{\bar\delta} \to \underline{\mathrm{Hom}}_{\delta}$
    est fidèle ;
  \item[(ii)] \struck{\ill{}} la restriction de $\varphi^{*}_C$ est surjective
    sur les $\underline{\mathrm{Hom}}$ ;
  \item[(ii$'$)] \struck{\ill{}} $\underline{\mathrm{Hom}}_{\delta}$ est dans
    l'image essentielle de $\underline{\mathrm{Hom}}_{\bar\delta}$.
\end{enumerate}

\textbf{N.B.}\quad (i) \struck{et (ii)} est dans le contexte \add{d'une
condition} \struck{d'} d'existence de flèches (\ill{} explicites de flèches) dans
$C$, \struck{et} (ii$'$) dans le contexte d'existence d'objets et de diagrammes
\ill{}. Ainsi (ii$'$) \struck{(ii)} \struck{par} \ill{} des conditions de type
(i$'$).

\ill{} généralisation de a) et b) [\struck{(ii)} \ill{} peut-on \add{ainsi}
exprimer \ill{} comme condition de type (ii$'$) ?]

Donc finalement on considère deux types de conditions permettant d'exprimer les
autres :
\begin{enumerate}
  \item[$1^{\circ}$)] \struck{\ill{}}
    $\mathrm{Image}\bigl(\varphi^{*}_C \mid \underline{\mathrm{Hom}}_{\bar\delta}\bigr)
    \subset \underline{\mathrm{Hom}}_{\delta}$ : impliquent conditions
    d'exactitude \struck{habituelles} ;
  \item[$2^{\circ}$)]
    $\mathrm{Image}\bigl(\varphi^{*}_C \mid \underline{\mathrm{Hom}}_{\bar\delta}\bigr)
    \supset \underline{\mathrm{Hom}}_{\delta}$ : impliquent \add{conditions d'}
    existence (y compris conditionnelle) des limites.
\end{enumerate}
Vérifier que cela permet d'exprimer l'égalité \ill{} resp. la surjectivité de
$\varphi^{*}_C \mid \underline{\mathrm{Hom}}_{\bar\delta}$ sur les
$\underline{\mathrm{Hom}}$ \ill{}

$2^{\circ}$) \emph{Conditions d'exactitude sur \struck{les} foncteurs
$F : C \to C'$.}

On a déjà vu comment on peut en exprimer avec des paquets de
\struck{diagrammes} \add{cônes} de $C$. Mais on veut des conditions
d'exactitude \ill{} pour tout $C$, sans donner de \struck{\ill{}}
supplémentaires sur $C$ — \ill{} : \add{un $C$} \ill{} paquets de
\struck{diagrammes} \add{cônes} de $C$. Une façon de le faire est de partir des
$\Theta$, muni de $\delta = (\delta',\delta'')$ et
$\lambda = (\lambda',\lambda'')$ \add{(paquets de cônes)}, \ill{} les
$u : \Theta \to C$ qui sont $\delta$-exacts, et de demander que $F$ soit
$u(\lambda)$-exact pour tout tel $u$ \struck{\ill{}}. On peut prendre un \ill{}
de \struck{\ill{}} $(\Theta_\alpha, \delta_\alpha, \lambda_\alpha) = T_\alpha$,
et $T = (T_\alpha)$. On dirait que \struck{\ill{}} la $T$-exactitude est
proprement \ill{}

\page{12}
Vu la composition \add{de foncteurs} \struck{entre} catégories \ill{}
satisfaisant aux mêmes conditions \ill{} demande de conditions d'exactitude.
Pour ceci, il est \struck{bon} naturel d'exiger que \struck{les foncteurs}
\ill{} $\lambda' \supset \delta'$, $\lambda'' \supset \delta''$.

\begin{tikzcd}[column sep=small, row sep=small, nodes={font=\scriptsize}]
C \arrow[r, "F"] & C' \arrow[r, "G"] & C'' \\
\Theta \arrow[u, "u"] & J' \arrow[l, "g'"'] & \\
\bar{I}' \arrow[u, "\xi'"] & \bar{I}'' \arrow[ul, "\xi''"'] & J'' \arrow[ull, "g''"']
\end{tikzcd}

\note{les deux $\bar{I}$ portent le paquet $\delta = (\delta',\delta'')$, les deux $J$ le paquet $\lambda = (\lambda',\lambda'')$}

Si on veut que les foncteurs identiques satisfassent les conditions, il est
naturel d'exiger qu'on se borne aux $C$ qui satisfont les conditions
d'exactitude \ill{} que $\forall\, \Theta \to C$ $\delta$-exact est
$\lambda$-exact. Mais alors la condition sur $F$ devient que
$\forall\, \Theta \xrightarrow{\ u\ } C$ $\lambda$-exact, $F \circ u$ est
$\lambda$-exact.

\struck{On a ainsi donné \ill{} et une \ill{} de définir :
$1^{\circ}$) \ill{} \add{une classe $\Sigma_0$} de parties \add{$\Delta_0$} de
$\mathrm{Ob}(\mathrm{Cat})$, chacune stable par l'équivalence des catégories
($C \in \Delta_0$, et $C' \in \mathrm{Ob}\,\mathrm{Cat}$, $C' \simeq C$, alors
$C' \in \Delta_0$) ; $\Sigma$ est stable par \ill{} construction.
$2^{\circ}$) Une classe $\Sigma_1$ de parties $\Delta_1$ de
$\mathrm{Fl}(\mathrm{Cat})$, chaque $\Delta_1$ stable par isomorphisme
($C \rightrightarrows C'$ par $u, v$, et $u \simeq v$, alors
$u \in \Delta_1 \Rightarrow v \in \Delta_1$), compatible avec l'identité, les
équivalences et la composition.}
\note{ce bloc est biffé d'un long trait diagonal}

\begin{enumerate}
  \item[$(1)$] Donnée d'un paquet \add{$P_1$} \add{de} catégories $\Theta$, avec
    $\lambda = (\lambda', \lambda'')$. Je dirai que \struck{\ill{}}
    $(\Theta, \lambda)$ \struck{$\Delta_\rho \subset (\mathrm{Cat})$ : les
    catégories \ill{}} $\Sigma_\rho \subset \mathrm{Fl}(\mathrm{Cat})$ : un
    foncteur $F : C \to C'$ est $\rho$-exact \add{$(\Theta,\lambda)$ dans $\rho$
    et $\forall$} si pour \struck{\ill{}} \add{tout foncteur}
    $u : \Theta \to C$ qui est $\lambda$-exact, $F \circ u$ est $\lambda$-exact.
\end{enumerate}
Alors
\begin{enumerate}
  \item[a)] $\Sigma_\rho$ est stable par isom. (i.e. si $F, G : C \to C'$, et
    $F \simeq G$, alors $F \in \Sigma_\rho \Rightarrow G \in \Sigma_\rho$) ;
  \item[b)] $\Sigma_\rho$ contient les équivalences de catégories ;
  \item[c)] $\Sigma_\rho$ stable par composition ;
  \item[d)] Une intersection des parties $\Sigma_\rho$ de
    $\mathrm{Fl}\,\mathrm{Cat}$ est de \add{ce} même type.
\end{enumerate}

\begin{enumerate}
  \item[$(2)$] Donnée d'un \add{couple $(\sigma_1, \sigma_2)$}
    \struck{de foncteurs} \add{de paquets de} foncteurs
    $R \xrightarrow{\ \varphi\ } \bar{R}$, avec \struck{\ill{}} $R$ muni de
    $\delta = (\delta', \delta'')$, $\bar{R}$ de
    $\bar\delta = (\bar\delta', \bar\delta'')$,
    \struck{\ill{}} \ill{} définit $\Delta_\sigma \subset \mathrm{Ob}\,\mathrm{Cat}$
    par :
\end{enumerate}

\page{13}
$C \in \Delta_\sigma$ ssi
\begin{enumerate}
  \item[a)] si $(R, \delta\,;\ \bar{R}, \bar\delta,\ \varphi : R \to \bar{R})
    \in \sigma_1$, alors
    \[
      \mathrm{Im}\bigl(\varphi^{*}_C \mid \underline{\mathrm{Hom}}_{\bar\delta}(\bar{R}, C)\bigr)
      \ \subset\ \underline{\mathrm{Hom}}_{\delta}(R, C) ;
    \]
  \item[b)] si $(\ \cdots\ ) \in \sigma_2$, alors on a l'inclusion inverse
    $\supset$.
\end{enumerate}
Alors
\begin{enumerate}
  \item[a)] $\Delta_\sigma$ est stable par équivalences de catégories ;
  \item[b)] Une intersection de \struck{catégories} \ill{} de la forme
    $\Delta_\sigma$ est de la même forme.
\end{enumerate}

\begin{enumerate}
  \item[$(3)$] $\Delta_{\sigma,\rho}$ est la \struck{\ill{}} sous-$2$-catégorie
    de $\mathrm{Cat}$ dont les objets sont les $\Delta_\sigma$ \ill{}
\end{enumerate}

\textbf{Remarque}\quad Soient $\Delta \subset \Delta'$ de la \ill{} ; pour
$C \in \mathrm{Ob}\,\Delta$, \add{un} $R \xrightarrow{\ \varphi\ } \bar{R}$ dans
$\Delta'$, et considérons
\[
  \underline{\mathrm{Hom}}_{\Delta'}(\bar{R}, C)
  \ \xrightarrow{\ \varphi^{*}_{\Delta',C} \,=\, \Psi\ }\
  \underline{\mathrm{Hom}}_{\Delta'}(R, C).
\]
Considérons l'une des conditions suivantes (sur $C$ \uncertain{brevité} de
$\mathrm{Ob}\,\Delta$) :
\begin{enumerate}
  \item[a)] $\Psi$ est fidèle ;
  \item[b)] $\Psi$ est surjectif sur les $\underline{\mathrm{Hom}}$ ;
  \item[c)] $\Psi$ est ess. surjectif sur objets.
\end{enumerate}
\note{à droite, deux accolades de sa main : a) et b) donnent « $\varphi$ pleinement fidèle », les trois ensemble « $\varphi$ est une équivalence »}

Soit $\Delta_0 \subset \Delta$ la sous-$2$-catégorie $2$-pleine définie par
l'une des trois conditions précédentes. Alors $\Delta_0$ est-elle distinguée ?

Soit \struck{$\Delta$} $\delta$ $(\bar\delta)$ l'un des \struck{diagrammes}
\add{cônes exacts} dans $R$ $(\bar{R})$ qui sont images des \struck{\ill{}}
\ill{} comme $\Theta$ pour $\Theta \to R$ ($\Theta \to \bar{R}$)
$\lambda$-exacts \struck{\ill{}}
\struck{distinctes des conditions dites} $(\Theta, \lambda) \in \rho'$. Donc
$\underline{\mathrm{Hom}}_{\Delta'}(\bar{R}, C) = \underline{\mathrm{Hom}}_{\bar\delta}(\bar{R}, C)$,
$\underline{\mathrm{Hom}}_{\Delta'}(R, C) = \underline{\mathrm{Hom}}_{\delta}(R, C)$.

\page{14}
\begin{enumerate}
  \item[(a)] $D = (D', D'')$ ens. de \struck{\ill{}} catégories.
  \item[(b)] $\sigma = (\sigma_1, \sigma_2)$, $\sigma_i$ $(i = 1,2)$ ens. de
    \add{quadruplets} $(R, \delta,\ \bar{R}, \bar\delta,\ \varphi : R \to \bar{R})$,
    $R$, $\bar{R}$ cat., $\varphi$ foncteur, i.e.
    $\delta = (\delta',\delta'')$ \struck{paires} \add{ens.} de cônes proj. et
    ind. de $R$ \add{$\bar{R}$} de types $\in D$,
    $\bar\delta = (\bar\delta', \bar\delta'')$ de même.
  \item[(c)] \struck{$\rho$ = paquet de $(\Theta, \lambda)$, avec $\Theta$ cat.,
    $\lambda = (\lambda', \lambda'')$ ens. de cônes projectifs et ind. de types
    $\in D$, et un foncteur $\rho$ contenant \ill{} au moins des $(R, \delta)$ et
    $(\bar{R}, \bar\delta)$ provenant de $\sigma$.}
\end{enumerate}

\page{15}
\note{feuillet biffé de deux grands traits en croix : une rédaction antérieure du théorème du feuillet 3}

\struck{On \ill{} \add{les} conditions découlant facilement de la suivante :
tout diagramme $\delta$ de type $\Theta \in d'$ ($\Theta \in d''$) dans $D$
parties d'un diagramme $\delta_j$ dans un des $C_j$ $(j < i)$, et que l'un des
indices $j < i$ \add{est} grand devant $\mathrm{card}\,\mathrm{Fl}\,\Theta$ : on
voit alors que la $\varinjlim$ des $D$ dans $C_j$ donne une limite de \ill{} des
$D$ [ou les $\varinjlim$ de type $\Theta$ commutent aux $\varinjlim$ grandes
devant $\mathrm{card}\,\mathrm{Fl}\,\Theta$].
b) $\alpha_i : C \to C_i$ est $\delta$-exact, i.e. si
$(x \xrightarrow{p_\lambda} x_\lambda)$ est dans $\delta'$, alors
$(\alpha_i(x) \xrightarrow{\alpha_i(p_\lambda)} \alpha_i(x_\lambda))$ est un cône
proj. exact. On est visiblement amené à la description des flèches de \ill{}}

\marginal{\ill{} comme le Pb 1, \ill{} est le Pb 2}

\textbf{Th.}\quad Soient $d'$, $d''$ des petits ensembles de diagrammes,
$\Delta = \Delta_{(d',d'')}$, $C$ une petite catégorie, munie d'un
$(\delta', \delta'')$ de cônes projectifs et inductifs. \struck{Alors} il existe
\struck{une $\Delta$-enveloppe de $(C, \delta', \delta'') = (\bar{C}, \ldots)$,
i.e. $i$ $(\delta',\delta'')$-exact, et universel vis-à-vis de cette propriété,
i.e. équivalence
$\underline{\mathrm{Hom}}_{\Delta, (\delta',\delta'')}(\bar{C}, \ldots) \to \underline{\mathrm{Hom}} \ldots$}
\struck{Soit $\rho$ le foncteur $\Delta_{(\emptyset, d'')} \to \ldots$,
$\sigma$ de même $\Delta_{(d', \emptyset)}$} — i.e. un foncteur
$v : C \to C'$ avec $C' \in \mathrm{Ob}\,\Delta$,
\struck{et telle \ill{} est petite}, vérifiant la propriété \add{qui + universel}
[i.e. $\forall\, E \in \mathrm{Ob}\,\Delta$,
\[
  \underline{\mathrm{Hom}}_{\Delta,\, (i(\delta'), i(\delta''))}(C', E)
  \;\xrightarrow{\ \approx\ }\;
  \underline{\mathrm{Hom}}_{\Delta,\, (\delta',\delta'')}(C, E)\ ].
\]

On construit une suite transfinie de catégories $C_i$ sous $C$, chacune munie
d'un ens. $\delta_i = (\delta'_i, \delta''_i)$ de cônes proj. et inductifs. On
pose $C_0 = C$, \struck{$\delta'_0, \delta''_0 = \ldots$}
$(\delta'_0, \delta''_0) = (\delta', \delta'')$ ;
$(C_{i+1}, \delta_{i+1})$ déduit de $(C_i, \delta_i)$ ainsi :

\page{17}
\section*{Topos et structures algébriques}

\begin{enumerate}
  \item[$1$)] Topos et structures algébriques \ill{}
\end{enumerate}
$\Delta_0$ = cat. avec $\varprojlim$ finies ; $\Delta$ = bons topos (conditions
$\supset$ b) c) de Giraud), foncteurs commutant aux $\varprojlim$, et exacts à g.
$T_0$ th. algébrique rel. $\Delta_0$, \struck{défini} \add{représentable} par une
petite $C \in \mathrm{Ob}\,\Delta_0$ \add{définissable par $\varprojlim$ finies} ;
$T_0(E) \simeq \underline{\mathrm{Hom}}_{\Delta_0}(C, E)$ si
$E \in \mathrm{Ob}\,\Delta_0$. Si $E \in \mathrm{Ob}\,\Delta$, on sait qu'on
\ill{}
\[
  T_0(E) = \underline{\mathrm{Hom}}_{\Delta_0}(C, E)
  \;\xrightarrow{\ \approx\ }\;
  \underline{\mathrm{Hom}}_{\Delta}(\hat{C}, E)
  \quad
  \bigl[\ \xrightarrow{\ \approx\ } \underline{\mathrm{Hom}}_{\mathrm{top}}(E, \hat{C})^{\circ}
  \ \text{si $E$ un topos}\ \bigr].
\]
Ainsi \add{le topos $\hat{C}$ est} \struck{\ill{}} la \struck{topos}
$\Delta_0 \to \Delta$ enveloppe de $C$.

\begin{enumerate}
  \item[$2$)] $\Delta_0$-\emph{Enveloppe d'une \add{petite} catégorie $I$}
    \quad (Pb. 1)
\end{enumerate}
On prend la $\Delta_0$-enveloppe
$\rho_{\Delta_0}(I) \subset \underline{\mathrm{Hom}}(I, \mathrm{Ens})^{\circ}$
(sous-cat. pleine des foncteurs « de présentation finie » — sous-catégorie
engendrée par $I$ \ill{} $\underline{\varprojlim}$ finies — on peut le former en
rajoutant les $\underline{\varprojlim}$ finies des diagrammes dans $I$, puis
prenant \ill{}). Puis on prend la $\Delta_0 \to \Delta$-enveloppe de
$\rho_{\Delta_0}(I)$, savoir $\widehat{\rho_{\Delta_0}(I)}$.

\begin{enumerate}
  \item[$3$)] \emph{Sous-théorie d'une $\Delta$-théorie} (représentable par un
    topos) \emph{et sous-topos} \quad (Pb 2.)
\end{enumerate}
Soit $R$ un topos, $\Sigma$ un ensemble de flèches de $R$. On considère
\[
  E \;\longmapsto\; \underline{\mathrm{Hom}}_{\Delta, \Sigma}(R, E),
  \qquad E \in \mathrm{Ob}\,\Delta .
\]
On \uncertain{satisfait} $\Sigma$ par ST1, ST2, ST3 (SGA 4 I 9.1.1 c)), d'où
$\Sigma'$, et on trouve $R\Sigma'^{-1}$, qui est un topos \add{$R'$}, et on a
\[
  \underline{\mathrm{Hom}}_{\Delta, \Sigma}(R, E)
  \;\xrightarrow{\ \approx\ }\;
  \underline{\mathrm{Hom}}_{\Delta}(R', E).
\]
On sait que $R'$ est un sous-topos de $R$, et tout sous-topos de $R$ s'obtient
ainsi (loc. cit.). Si $R \simeq \hat{C}$, $C$ petite \ill{}

\marginal{ST1\quad Si $w = vu$, \struck{\ill{}} et \struck{\ill{}} des flèches
\struck{\ill{}} \struck{sont} dans $\Sigma$, \struck{\ill{}} le troisième l'est.
Contravariance}
\marginal{ST2\quad Stabilité par changement de base (\struck{\ill{}} par
$\underline{\varprojlim}$ finies)}
\marginal{ST3\quad Nature locale ($\iff$ stabilité par $\underline{\varprojlim}$)}

\page{18}
ils correspondent aux topologies $T'$ \add{sur $C$} \struck{plus} plus fines que
la top. donnée, par $R' \longleftrightarrow (C, T')$.

Remarquons que pour un sous-topos donné \add{$R'$}, on peut toujours le décrire à
l'aide d'un petit ens. $\Sigma \subset \mathrm{Fl}\,R$. En effet, sur
$C \subset R$ \add{petite et} \struck{génératrice} \ill{}. Alors les flèches
$\in \Delta$, $R \to \bar{R}$, qui se factorisent par $R'$ sont celles qui
transforment familles \struck{couvrantes} $S_i \to S$ \struck{\ill{}}
\add{dans $C$, couvrantes pour la topologie $T'$ définissant $R'$} en familles
couvrantes [\ \ill{} \add{une} telle que $\forall\, S \in \mathrm{Ob}\,C$ et
$S' \to S$ monomorphisme de \struck{\ill{}} $R$, $S' \to S$ soit transformé en
mono ].

\textbf{Corrigeons}\quad Supposons $R$ « algébrique » i.e. $T_R$
\add{théorie} « algébrique » ; alors $T_{R'} \subset T_R$
(\struck{sous-catégorie} \ill{} théorie « pleine ») est algébrique, déduite de
$T_R$ par un petit ensemble d'axiomes.

D'ailleurs, si on a deux théories représentables $T_R$, $T_{R'}$ ($R$, $R'$ des
topos) et un foncteur $T_{R'} \to T_R$ qui est un plongement de théories, i.e.
tel que \struck{\ill{}} $T_{R'}(E) \to T_R(E)$ pl. fid. pour
$\forall\, E \in \mathrm{Ob}\,\Delta$, ici \ill{} \ill{} dans le \ill{}
correspondant $R \to R'$ de $\Delta$ est telle que
$\forall\, E \in \mathrm{Ob}\,\Delta$,
$\underline{\mathrm{Hom}}_{\Delta}(R', E) \to \underline{\mathrm{Hom}}_{\Delta}(R, E)$
est pl. fidèle. On \ill{} caractérise les morphismes de topos $R' \to R$ qui sont
des plongements.

\begin{enumerate}
  \item[$4$)] \textbf{Théorème}
  \item[a)] Un $R \in \mathrm{Ob}\,\Delta$ est \ill{} algébrique, i.e.
    représente une théorie algébrique, ssi $R$ est un topos, i.e. admet une
    petite sous-catégorie génératrice.
  \item[b)] Soit $I$ un ens. discret, $R \in \mathrm{Ob}\,\Delta$,
    $I \xrightarrow{\ b\ } R$ un foncteur. Alors $b$ est \emph{basique}
    (i.e. \ldots) \struck{ssi} \struck{disque} \struck{$b(I)$ est génératrice}
    ssi les sous-catégories pleines de $R$ engendrées par produits finis
    d'objets de $b(I)$ sont strictement génératrices.
\end{enumerate}

\textbf{Dém.}\quad \struck{\ill{}} Évidemment b) $\Rightarrow$ a), donc
\struck{\ill{}} précisons \struck{\ill{}} b).

\emph{Suffisance}\quad On va construire la théorie alg. \add{basique}
représentée par $R$, $b$ en plusieurs pas.

\page{19}
\begin{enumerate}
  \item[$1^{\circ}$)] \struck{Théorie} Structures vides sur objets de base
    $(X_i)_{i \in I}$. Représentables par un topos $R_0 =$ (Pb 1 ; on a
    $R_0$ = ens. \struck{\ill{}} $I$-multisimpliciaux \add{$\omega$}).
    $J$ est génératrice dans $R$ :
    $I \to J_0 \to R_0 = \hat{J_0} \to R$.
    \struck{\ill{}} un objet de $R_0$ provenant de \ill{}
  \item[$2^{\circ}$)] On rajoute des flèches \struck{\ill{}} $J_1$ et
    $R_1 = \rho_\Delta(J_1)$ ; $J_0$ (celles de $R$), d'où $R_0$ \struck{\ill{}}
  \item[$3^{\circ}$)] on met des axiomes d'égalité sur des composés (en vertu de
    $R$), d'où $J \to R_0 \to R_1 \to R_2 \to R$ et
    \[
      T_{R_2}(E) = \underline{\mathrm{Hom}}_{\Delta}(R_2, E)
      \;\xrightarrow{\ \approx\ }\; \underline{\mathrm{Hom}}(J_1, R)
      \;\longleftrightarrow\; \underline{\mathrm{Hom}}_{\Delta}(R, E) = T_R(E)
    \]
    où $J_1$ est la sous-catégorie pleine de $R$ définie par $b : I \to R$, dans
    la petite catégorie $J_1$ (cf. Pb. 1). $R_2$ est la $\Delta$-enveloppe de la
    \ill{}. \add{Comme} $T_R \to T_{R_2}$ est pl. fidèle, on sait qu'on a fini
    par \ill{}
\end{enumerate}
\marginal{catégorie \ill{} $\rho(J_0) = \rho(J'_1)$, $J'_1$ analogue engendré par $J_0$ et les flèches données}

\struck{\ill{}} Maintenant \struck{\ill{}} \add{les foncteurs} \ill{} ne
prennent que \ill{} \struck{\ill{}} que la structure. [on peut

\begin{enumerate}
  \item[$4^{\circ}$)] On met d'autres axiomes \struck{\ill{}} explicites,
    lesquels : a) foncteurs \underline{continus} $I_1 \to E$, i.e.
    \struck{\ill{}} qui transforment certaines \struck{systèmes}
    \add{(aux cônes exacts dans $R$)} cônes inductifs de $I_1$
    \struck{sans cônes} exacts ; \struck{b) foncteurs} en \ill{} un des
    foncteurs \ill{} qui ne peuvent \ill{} correspondent aux
    \underline{$\varprojlim$} $R \to E$, plus la condition de transformer
    certains \struck{diagrammes de} cônes projectifs \ill{} de $R'_2$, le type
    classifiant par \ill{} cônes exacts \ldots
\end{enumerate}
On trouve $I \to R_0 \to R_1 \to R_2 \to R_3 \simeq R$.

\marginal{$\hat{J_1} \to R$ est un foncteur de \underline{localisation}}

\textbf{N.B.}\quad Comme on a rajouté \ill{} $3^{\circ}$ et $4^{\circ}$ des
axiomes, on peut les condenser en un seul pas, \ill{} qu'en fait
$R_1 \to R_2$ est surjectif \ill{} objets et flèches (cf. Pb 2). Donc on peut
faire \struck{\ill{}} par les trois pas minimum : objets, \struck{\ill{}}
flèches \add{aux produits finis}, axiomes sur les flèches ; exprimables en
termes \struck{explicites} d'insertion de flèches entre \ill{} en produits)
\struck{\ill{}} \ill{} telles familles qu'aux flèches introduites dans les
données.

\page{20}
\struck{\ill{}}

\textbf{Re N.B.}\quad Si on veut seulement \ill{} \add{une théorie} pour \ill{}
un topos $R$ soit « algébrique », c'est encore plus court : on choisit $C$ petite
avec $\underline{\varprojlim}$ finies, telle que $R \simeq \hat{C}$, et
\struck{on voit que} par la structure algébrique $T^{\Delta_0}_C$ — sa
restriction \ill{}, $\Delta$ est \ill{} algébrique, elle est représentée par
$\hat{C}$, \struck{sauf} \add{$n = 1$} — puis on rajoute des axiomes exprimant
la \underline{continuité} d'un foncteur $C \to E$ (commutant déjà aux
$\underline{\varprojlim}$ finies) par le fait de transformer les familles
couvrantes en \ill{} familles couvrantes \ldots \struck{\ill{}}

\struck{Nécessité}\quad Il faut prouver que :

\begin{enumerate}
  \item[a)] Structure vide sur un ens. de base $I$ \add{par la construction
    transfinie}. Représentée par $\hat{\tilde{I}}$ (cf. Pb 1), où
    $\tilde{I} = \rho_{\Delta_0}(I)$.
    \struck{$\tilde{I}$ $\Delta$-engendre} \textbf{N.B.} $\tilde{I}$
    $\Delta$-engendre $\hat{\tilde{I}}$ — deux pas : \ill{} rajouter aux
    produits, puis aux $\underline{\varprojlim}$ quelc.
  \item[b)] \struck{\ill{}} Mettre axiomes sur théorie algébrique $T_R$, $R$ un
    topos, i.e. \ill{} trouver $R'$ et $R \to R'$ en surjectif (et \ill{} plus),
    ainsi surjectif sur les flèches \ill{} des Pb 2 traité. On trouve \ill{}
  \item[c)] Reporter flèches (petit ensemble \add{$H$}) : $R$ \ill{} théorie
    algébrique $T_R$, $R$ déjà topos. \struck{Soit $C$ l'ens. des} Soit
    \struck{$C \subset R$ la petite sous-catégorie} \add{$C$ $\mathrm{Ob}\,R$}
    \ill{} d'une des sources et buts des \ill{} flèches, $C$ la catégorie libre
    engendrée par les objets et les flèches \ill{} $H$ ; $\rho_\Delta(C)$ la
    catégorie \ill{} ne modifiant \ill{} dans la théorie $T'$ est le produit
    $2$-fibré
    \[
      T' \simeq T_R \times^{2}_{T_I} T_C
    \]
    \struck{\ill{}} où
    $T_C(E) \simeq \underline{\mathrm{Hom}}_{\mathrm{Cat}}(C, E)$, et $T_I(E)$
    se définit de même, $I$ considéré comme catégorie discrète. Or (Pb 1)
    $T_C \simeq T_{R_1}$, $T_I \simeq T_S$, $R$ et $S$ deux topos, donc il faut
    prouver l'existence \ill{}
\end{enumerate}

\end{document}
